\documentclass[../Main.tex]{subfiles}

\begin{document}
In this chapter let $A \in M_{n \times n}(\F)$.
\section{Defining the Determinant}
\begin{definition}{Volume form}
    A function $F : M_{n \times n}(\F) \mapsto \F$ is a \underline{volume form} if it is:
    \begin{enumerate}
        \item \textbf{alternating}: if $\vec{c_{i}}, \vec{c_{j}}$ are equal columns of $A$ with $i \neq j$, then $F(A) = 0$;
        \item \textbf{multilinear in columns}: For all $i \in \{1, \cdots, n\}$ and $\vec{v_j} \in \F^n$, the restricted function:
            \begin{align*}
                \F^n &\mapsto \F\\
                \vec{v} &\mapsto F(V)
            \end{align*}
            is linear, where $V$ is the matrix with columns $\vec{v_j}$ and $\vec{v}$ as the $i$th column.
    \end{enumerate}
\end{definition}
\begin{propositions}{
        Let $F$ be an $n$-dimensional volume form.
        \label{propsVolFormProps}
    }
    \item If $A$ has a zero column then $F(A) = 0$;
    \item $F(AT_{ij}) = -F(A)$;
    \item $F(AM_{i, \lambda}) = \lambda F(A)$;
    \item $F(AC_{i, j, \lambda}) = F(A)$.
\end{propositions}
\begin{proof}
    The proofs a simple from the definitions of a volume form. It is helpful to define $\vec{c_{i}}$ as the columns and $f_i(\vec{v})$ as the function evaluating $F$ on $A$ except column $i$ is set to $\vec{v}$.
\end{proof}
\begin{theorem}[Determinant uniqueness theorem]
    There is a unique volume form $F : M_{n \times n}(\F) \mapsto \F$ satisfying non-triviality, $F(I_n) = 1$.
    \label{thmDeterminant}
\end{theorem}
\begin{proof}
    First note that \thmref{propsVolFormProps} gives us that $F(AE) = F(A) F(E)$ for elementary $E$. We can show uniqueness of $F$ by considering the RRE form of $A$ and expanding and arguing uniqueness of RRE form.

    To show existence of $F$, define:
    \begin{equation*}
        F(A) = \sum_{\sigma \in S_n} \sgn(\sigma) \prod_{i=1}^n a_{\sigma(i) i}
    \end{equation*}
    and check that it satisfies the required properties. Show the alternating form by considering a transposition $(i~~j)$.
\end{proof}
\begin{definition}{Determinant}
    The \underline{determinant} of $A$ is $\det(A)$ or $|A|$ and is the function defined in \thmref{thmDeterminant}.
\end{definition}
\begin{corollary}
    $\det(A) \neq 0$ if and only if $A$ is invertible.
    \label{corDetInvert}
\end{corollary}
\begin{proof}
    If $A$ is invertible then it is a product of elementary matrices which have non-zero determinants. Use reversible logic.
\end{proof}
\section{Properties of the Determinant}
\begin{lemma}
    Determinant is unaffected by transpose:
    \begin{equation*}
        \det(A) = \det(A^T)
    \end{equation*}
    \label{lemDetTranspose}
\end{lemma}
\begin{proof}
    Use the definition of the determinant directly.
\end{proof}
\begin{proposition}
    For $A, B \in M_{n \times n}(\F)$,
    \begin{equation}
        \det(AB) = \det(A) \det(B)
        \label{eqnDetProduct}
    \end{equation}
    \label{propDetProduct}
\end{proposition}
\begin{proof}
    If $A$ or $B$ does not have full rank, then by \thmref{lemRankComposition} neither does $AB$ and so both sides of \eqnref{eqnDetProduct} are zero.

    Otherwise, write $A$ and $B$ as products of elementary matrices for which \eqnref{eqnDetProduct} already holds.
\end{proof}
\begin{corollary}
    If $A$ is invertible then:
    \begin{equation*}
        \det(A^{-1}) = \frac{1}{\det(A)}
    \end{equation*}
    \label{corDetInverse}
\end{corollary}
\begin{proof}
    \begin{align*}
        1 &= \det(I_n) = \det(AA^{-1}) \\
        &= \det(A) \det(A^{-1})
    \end{align*}
\end{proof}
\begin{corollary}
    Similar matrices have the same determinant
    \label{corDetSimilar}
\end{corollary}
\begin{proof}
    Trivial using \thmref{propDetProduct}.
\end{proof}
\begin{definition}{Determinant (linear map)}
    For a linear map $\alpha \in \L(V, V)$ represented by matrix $A$ in some basis, the \underline{determinant} of $\alpha$ is:
    \begin{equation*}
        \det(\alpha) = \det(A)
    \end{equation*}
    and the choice of basis does not matter by \thmref{corDetSimilar}.
\end{definition}
\begin{definition}{Singularity (linear map)}
    A linear map is \underline{singular} if its determinant is zero.
\end{definition}
\begin{proposition}
    Let $\alpha \in \L(V, V)$. Then the following are equivalent:
    \begin{enumerate}
        \item $\rk(\alpha) = \dim(V)$;
        \item $\alpha$ is invertible;
        \item $\alpha$ is non-singular.
    \end{enumerate}
    \label{propLinMapDetEquivs}
\end{proposition}
\begin{proof}
    Easy by considering any matrix of $\alpha$.
\end{proof}
\begin{proposition}[Block-triangular determinant]
    Let $A \in M_{k \times k}(\F), B \in M_{l \times l}(\F)$. Define:
    \begin{equation*}
        C = 
        \begin{pmatrix}
            A & (\star) \\
            0 & B
        \end{pmatrix}
    \end{equation*}
    where the star represents any matrix. Then $\det(C) = \det(A) \det(B)$.
    \label{thmBlockTriDet}
\end{proposition}
\begin{proof}
    Consider the determinant definition and note that the only elements of $S_n$ that contribute are those that preserve $\{1, \cdots, k\}$. Therefore, it must be made up of $\sigma$ and $\tau$ which commute and preserve $\{1, \cdots, k\}$ and $\{k+1, \cdots, k + l\}$ respectively. Then factorise the determinant expression using these.
\end{proof}
We can also inductively apply this to any upper-triangular block matrix.
\begin{definition}{Minor matrix}
    The \underline{$(i, j)$ minor matrix} is:
    \begin{equation*}
        \mnr{A}{i}{j} \in M_{(n - 1) \times (n - 1)}(\F)
    \end{equation*}
    obtained by deleting the $i$th row and $j$th column.
\end{definition}
\begin{theorem}[Laplace Expansion]
    For any fixed $i \in \{1, \cdots, n\}$,
    \begin{equation*}
        \det(A) = \sum_{j=1}^{n} (-1)^{i+j} a_{ij} \det(\mnr{A}{i}{j})
    \end{equation*}
    and similar for expanding about a column.
    \label{thmLaplaceExpansion}
\end{theorem}
\begin{proof}
    We only need to show the column expansion because we can apply this to the transpose by \thmref{lemDetTranspose}. Therefore fix $j$.

    If $A^{(i)}$ is the matrix with column $j$ set to $\vec{e_i}$, then:
    \begin{equation*}
        \det(A) = \sum_{i=1}^{n} a_{ij} \det(A^{(i)})
    \end{equation*}
    Use row and column operations to then get $A^{(i)}$ to have the form:
    \begin{equation*}
        \begin{pmatrix}
            1 & (\star) \\
            0 & \mnr{A}{i}{j}
        \end{pmatrix}
    \end{equation*}
    with some $-1$ pre-factors, simplify for the required form.
\end{proof}
\begin{definition}{Adjugate matrix}
    The \underline{adjugate matrix} $\adj(A)$ of $A$ is given by:
    \begin{equation*}
        \left(\adj(A)\right)_{ij} = (-1)^{i+j} \det(\mnr{A}{j}{i})
    \end{equation*}
\end{definition}
\begin{theorem}[Cramer's Rule]
    \begin{enumerate}
        \item $\adj(A) A = A \adj(A) = \det(A) I$
        \item If $A$ is invertible then:
            \begin{equation*}
                A^{-1} = \frac{\adj(A)}{\det(A)}
            \end{equation*}
    \end{enumerate}
    \label{thmCramer}
\end{theorem}
\begin{proof}
    2 follows from 1 by uniqueness of the inverse.

    For 1, first use \thmref{thmLaplaceExpansion} to expand around the $j$th column and show that the diagonal entries of $\adj(A) A$ are $\det(A)$.

    Then consider $A'$ which is the same as $A$ except has the $k$th column also in the $j$th position. This has determinant zero, and by expanding the determinant around the $j$th column again we get that the off-diagonal entries of $\adj(A) A$ are zero.

    For the other expression we expand around rows.
\end{proof}
\section{Trace}
\begin{definition}{Trace}
    The \underline{trace} of $A$ is given by:
    \begin{equation*}
        \tr(A) = \sum_{i=1}^{n} a_{ii}
    \end{equation*}
\end{definition}
Trace is a linear operation on matrices. 
\begin{proposition}
    For $A, B \in M_{n \times n}(\F)$,
    \begin{equation*}
        \tr(AB) = \tr(BA)
    \end{equation*}
    \label{propTraceCommute}
\end{proposition}
\begin{proof}
    Simply by expanding the required terms.
\end{proof}
Note that $\tr(AB) \neq \tr(A)\tr(B)$ in general.
\begin{corollary}
    Similar matrices have the same trace
    \label{corTraceSimilar}
\end{corollary}
\begin{proof}
    Consider $\tr(PAP^{-1})$ and cycle the terms in the trace for cancellation.
\end{proof}
\begin{definition}{Trace (linear map)}
    A linear map $\alpha : V \mapsto V$ has \underline{trace}:
    \begin{equation*}
        \tr(\alpha) = \tr([\alpha]_B^B)
    \end{equation*}
    where $B$ is any basis for $V$. This is independent of the chosen basis.
\end{definition}
\end{document}